\documentclass[10pt,a4paper]{amsart}
\usepackage[margin=19mm]{geometry}
\usepackage{amsmath,amssymb,mathtools}
\usepackage[scr=rsfs,cal=euler]{mathalfa}
\usepackage{xfrac}
\usepackage[unicode,hidelinks,bookmarks=false]{hyperref}
\newcommand{\ie}{i.e.\ }
\newcommand{\cf}{cf.\ }
\newcommand{\resp}{respectively\ }
\newcommand{\Subsection}{\subsection}
\providecommand{\nobreakdash}{\nobreak\hbox{-}}
\setlength{\emergencystretch}{2em}
\multlinegap0pt
\usepackage{xcolor}
\usepackage{amssymb}
\usepackage{enumerate}
\newtheorem{lemma}{Lemma}
\newcommand{\NN}{N}
\newcommand{\PP}{\mathcal P}
\newcommand{\RE}{\mathbb R}
\renewcommand{\d}{\mathrm{d}}
\DeclareMathOperator{\sgn}{sgn}
\newcommand{\x}{{\bf  x}}
\title{The Born-Oppenheimer approximation for a 1D 2+1 particle system with zero-range interactions}
\author{Claudio Cacciapuoti, Andrea Posilicano, Hamidreza Saberbaghi}
\date{Source: arXiv:2506.21457v2 (CC BY 4.0)}
\begin{document}
\maketitle

\noindent\emph{Source and licence.} The Born-Oppenheimer approximation for a 1D 2+1 particle system with zero-range interactions. Version arXiv:2506.21457v2 (CC BY 4.0), distributed by arXiv under the Creative Commons Attribution 4.0 International licence, http://creativecommons.org/licenses/by/4.0/, as recorded in the arXiv licence metadata for this exact version and shown on the abstract page https://arxiv.org/abs/2506.21457v2. Full licence notice: \texttt{licenses/arxiv-2506.21457-CC-BY-4.0.txt}.\par\smallskip

\noindent\textit{Editorial scope.} Counted unit: the article's lemma l:PH (source0.tex lines 1433--1443). The article's complete original proof of it is delivered with it (source0.tex lines 1444--1565, one \texttt{proof} environment of 122 lines). No mathematics is written, repaired or re-derived: the statement and the proof are the article's own lines, in the article's own notation, and no retained line is substituted or re-flowed.  Retained uncounted as the source-local context the proof needs: lines 1366--1369; lines 1371--1373.  Nothing was omitted from any retained span, so the package carries no omission record.  No retained line contains a citation, so the wrapper carries no bibliography and no bibliography entry.  No retained line contains a figure environment, an image-inclusion command or drawing code, so no image is delivered and the package declares no asset dependency.  Editorial formatting only, in addition: an AMS \texttt{amsart} preamble that restates the article's own declarations and macros used by the retained lines, namely the environments \texttt{lemma} on separate counters, together with the standard packages those lines need.  The retained environments print in this document as Lemma 1 in order of appearance, whereas the article prints the same items with the numbering of its own full document.  The complete native source of the article as distributed in the arXiv e-print for this exact version is bundled unchanged as \texttt{sources/180564/source0.tex}.
\begin{equation}\label{psiBOx}
 \psi^{BO}_{x}(y) :=    N(x) 
\left(e^{-\sqrt{\lambda_{0}(x)}\,|x/2-y|} + e^{-\sqrt{\lambda_{0}(x)}\,|x/2+y|}\right) ,
\end{equation}

 \begin{equation}\label{Normalization_constant}
        N(x)  := \left(\frac{\sqrt{\lambda_{0}(x)}}{2\left(1+ e^{-\sqrt{\lambda_{0}(x)}\,|x|}\big(1+ \sqrt{\lambda_{0}(x)}\,|x|\big)\right)} \right)^{1/2} .
    \end{equation}

\begin{lemma}\label{l:PH}  i) $\mathcal P$ is a bounded operator in $H^1(\RE^2)$; 
ii) the commutator $[\partial_{x},\PP]:H^{1}(\RE^{2})\subset L^2(\RE^2)\to L^2(\RE^2)$ extends to a bounded operator in $L^{2}(\RE^{2})$ and 
\begin{equation*}%\label{commutator-bound}
\|\,[\partial_{x}, \PP]\,\| \leq \delta\, ,
\end{equation*}
with 
\begin{equation*}%\label{delta}
\delta := 2\left( \sup_{x\in\RE} \int |\partial_x \psi^{BO}(x,y)|^2 \d y \right)^{1/2}
\end{equation*}
and the bound $\delta \leq 4 |\alpha|$.
\end{lemma}

\begin{proof}
In the course of the proof, for the sake of brevity, we will often omit the dependence on $x$ in $\NN(x)$ and $\lambda_0(x)$. From the formula \eqref{psiBOx}, one gets
\[
\partial_y \psi^{BO}(x,y) =     {\NN}\sqrt{\lambda_0 } 
\left(\sgn(x/2-y)e^{-\sqrt{\lambda_0 }\,|x/2-y|} - \sgn(x/2+y)e^{-\sqrt{\lambda_0 }\,|x/2+y|}\right) ,
\]
and so
\[\begin{aligned}
\int_\RE \left|\partial_y\psi^{BO}(x,y)\right|^2 \d y  
= & 2 \NN^2 \sqrt{\lambda_0 }\left(1+e^{-\sqrt{\lambda_0 }\,|x|} -\sqrt{\lambda_0 }\,|x| e^{-\sqrt{\lambda_0 }\,|x|} \right) \\ =&  \lambda_0 \,\frac{1+e^{-\sqrt{\lambda_0 }\,|x|} -\sqrt{\lambda_0 }\,|x| e^{-\sqrt{\lambda_0 }\,|x|} }{1+e^{-\sqrt{\lambda_0 }\,|x|} +\sqrt{\lambda_0 }\,|x| e^{-\sqrt{\lambda_0 }\,|x|} } \\
\leq &\lambda_0  \leq\alpha^2\,.
\end{aligned}\]
Hence,
\begin{align*}
\int_{\RE^2}\left|\partial_y\mathcal{P} \phi(x,y)\right|^2 \d\x  = &\int_{\RE^2} \left|\partial_y\psi^{BO}(x,y)\right|^2 |f_\phi(x)|^2 \d x\,\d y  \\
  \leq & \left(\sup_{x \in \RE} \int_\RE \left|\partial_y\psi^{BO}(x,y)
\right|^2 dy \right)\|f_{\phi}\|^2_{L^2(\RE)}\\
\leq &\alpha^{2}\| \phi \|^2_{L^2(\RE^2)}\, .
\end{align*} 
Defining
\begin{equation*}%\label{tildef}
    \tilde f_\phi(x) :=  \int_{\RE} \big(\partial_x\psi^{BO}(x,y)\big) \phi(x,y) \d y \,,
    \end{equation*}
one has
     \[
\|  \tilde f_\phi\| \leq \frac{\delta}{2}\,\|\phi\|\,. 
    \]
    Since
    \begin{equation*}%\label{partialxPphi}
    \partial_x \PP\phi  = \big(\partial_x\psi^{BO}\big) f_\phi  +\psi^{BO}\tilde f_\phi +  \psi^{BO}f_{\partial_x\phi},
    \end{equation*}
    there holds 
  \[\begin{aligned}
    \|\partial_x \PP \phi\|  \leq & \|\big(\partial_x\psi^{BO}\big) f_\phi\|  +\|\psi^{BO}\tilde f_\phi \| +\|\psi^{BO} f_{\partial_x\phi} \|  \\ 
     \leq &
     \frac\delta2\,\|f_\phi\|  +\|\tilde f_\phi \| + \|\partial_x \phi \| \leq (1+\delta) \|\phi\|_{H^1(\RE^2)} . 
   \end{aligned}
   \]
    In a similar way, one obtains 
 \[
    [\partial_x, \PP]\phi  = \big(\partial_x\psi^{BO}\big) f_\phi  +\psi^{BO}\tilde f_\phi ,
    \]
    and   
    \[
    \|\,[\partial_x, \PP]\phi\|  
     \leq \delta\, \|\phi\|. 
    \]   
We are left to prove the bound on $\delta$.
    From the formula \eqref{psiBOx}, one can explicitly compute $\partial_x\psi^{BO}$:
\[
\begin{aligned}
\partial_x \psi^{BO}(x,y) =&
\underbrace{\frac{\NN'}{\NN} \psi^{BO}}_{\varphi^{BO}_1}
-
\underbrace{\frac{{\NN}\lambda_0'}{2\sqrt{\lambda_0}}  
\left(|x/2-y|e^{-\sqrt{\lambda_0}|x/2-y|} +|x/2+y|e^{-\sqrt{\lambda_0}|x/2+y|}\right)}_{\varphi^{BO}_2} 
\\
&-\underbrace{\frac{{\NN} \sqrt{\lambda_0}}{2} 
\left(\sgn(x/2-y)e^{-\sqrt{\lambda_0}|x/2-y|} + \sgn(x/2+y)e^{-\sqrt{\lambda_0}|x/2+y|}\right)}_{\varphi^{BO}_3}. 
\end{aligned}
\]
Computing the norm of $\varphi^{BO}_1$ one obtains
\[
\|\varphi^{BO}_1\|^{2} =\left( \frac{\NN'}{\NN}\right)^2.
\]
From Eq. \eqref{Normalization_constant}, by a straightforward calculation there follows 
    \begin{equation*}
            \frac{\NN'}{\NN} = \frac{(\sqrt{\lambda_0})^{\prime}}{2\sqrt{\lambda_0}} + \frac{\frac{(\sqrt{\lambda_0})^{\prime}}{\sqrt{\lambda_0}} \left(\lambda_0 x^2 + \frac{\lambda_0^{3/2}}{(\sqrt{\lambda_0})^{\prime}}x\right)}{2 \left(  \sqrt{\lambda_0}|x|+e^{\sqrt{\lambda_0}|x|}+1\right)} \, .
    \end{equation*}
We notice that 
\[
\lambda_0(x) = \frac{\alpha^2}{4} \nu^2(|\alpha| x/ 2) \qquad x\geq 0
\]
with 
\[
\nu(x) = \frac{W(xe^{-x})}{x}+1.
\]
Hence, 
\begin{equation*}%\label{ratio_bound}
    \left\lvert\frac{(\sqrt{\lambda_0})^{\prime}}{2\sqrt{\lambda_0}} \right\rvert = \frac{|\alpha|}{4}  \left|\frac{\nu'(ax/2)}{\nu(ax/2)}\right|.  
\end{equation*}
From the identity $W(x)e^{W(x)} = x$, $x\geq 0$, we infer 
\[
\nu(x) = e^{-x\nu(x)}+1.
\]
Hence, 
\[
\frac{\nu'(x)}{\nu(x)} = - \frac{1}{e^{x\nu(x)}+x}
\]
and $|\nu'(x)/\nu(x)| \leq \frac{1}{e^{x}+x} \leq 1$. We deduce the bound:  
\begin{equation} \label{sqrtlambdaprime}
    \left| \frac{(\sqrt{\lambda_0})^{\prime}}{2\sqrt{\lambda_0}} \right| \leq \frac{|\alpha|}{4}.
\end{equation}
Recalling also that $\sqrt{\lambda_0} \leq |\alpha|$, we infer
    \begin{equation*}%\label{111}
 \left|\frac{\frac{(\sqrt{\lambda_0})^{\prime}}{\sqrt{\lambda_0}} \left(\lambda_0 x^2 + \frac{\lambda_0^{3/2}}{(\sqrt{\lambda_0})^{\prime}}x\right)}{2 \left(  \sqrt{\lambda_0}|x|+e^{\sqrt{\lambda_0}|x|}+1\right)} 
    \right| 
    \leq \frac{\frac{|\alpha|}{2}\lambda_0 x^2+|\alpha|\sqrt{\lambda_0}|x|
    }{2 \left(  \sqrt{\lambda_0}|x|+e^{\sqrt{\lambda_0}|x|}+1\right)}\leq \frac{|\alpha|}{4}, 
        \end{equation*}
        where we used the  inequality $\frac{\frac{1}{2}s^2+s}{2 (s +e^{s}+1)} \leq \frac14$, $s\geq 0$. We deduce that   \[ \left| \frac{\NN'}{\NN}\right|  \leq \frac{|\alpha|}{2}\qquad \text{and} \qquad  
        \|\varphi^{BO}_1\|_{L^2(\RE,dy)}^2 =  \left( \frac{\NN'}{\NN}\right)^2  \leq \frac{\alpha^2}{4} .\]
Next we focus attention on $\varphi^{BO}_2$.
\[
\begin{aligned}
\|\varphi^{BO}_2\|^2 =& \frac{({\NN}\lambda_0')^2}{4(\lambda_0)^{5/2}}\left(1+\frac{(\sqrt{\lambda_0}|x|)^3}{3}e^{-\sqrt{\lambda_0}|x|} +(1+\sqrt{\lambda_0}|x|) e^{-\sqrt{\lambda_0}|x|} \right) \\
=&
\frac{(\lambda_0')^2}{2(\lambda_0)^{2}}\frac{1+\frac{(\sqrt{\lambda_0}|x|)^3}{3}e^{-\sqrt{\lambda_0}|x|} +(1+\sqrt{\lambda_0}|x|) e^{-\sqrt{\lambda_0}|x|}}{1+e^{-\sqrt{\lambda_0}|x|} +\sqrt{\lambda_0}|x| e^{-\sqrt{\lambda_0}|x|}} \leq \frac{(\lambda_0')^2}{(\lambda_0)^{2}} \leq\alpha^2,
\end{aligned}
\]
where in the latter step we used the bound  $|\lambda_0'/\lambda_0| \leq |\alpha|$ that can be easily obtained from the bound \eqref{sqrtlambdaprime}. 

We conclude with a bound on $\|\varphi^{BO}_3\|$.
\[
\begin{aligned}
\|\varphi^{BO}_3\|^2  =& \frac{{\NN}^2\sqrt{\lambda_0}}{2}\left(1-e^{-\sqrt{\lambda_0}|x|} +\sqrt{\lambda_0}|x| e^{-\sqrt{\lambda_0}|x|} \right) \\
=&
 \frac{\lambda_0}{4}\frac{1-e^{-\sqrt{\lambda_0}|x|} +\sqrt{\lambda_0}|x| e^{-\sqrt{\lambda_0}|x|} }{1+e^{-\sqrt{\lambda_0}|x|} +\sqrt{\lambda_0}|x| e^{-\sqrt{\lambda_0}|x|}} \leq  \frac{\lambda_0}{4} \leq \frac{\alpha^2}{4}.
\end{aligned}
\]
Summing up we obtain $\|\partial_x \psi^{BO}\|  \leq \sum_{j=1}^3 \|\varphi_j^{BO}\|  \leq 2|\alpha|$ and $\delta \leq 4|\alpha| $. 
\end{proof}

\end{document}
